% Candidate continuation fragment for gaga.tex covering GAGA source units
% gaga:no:5 through gaga:lemma:3. Insert this fragment immediately before
% the unique terminal \input{chapters} in the foundation candidate.
%
% Convention: throughout this fragment an algebraic variety over C means a
% reduced separated scheme of finite type over C, not necessarily irreducible.

\section{Analytification of algebraic varieties}
\label{section-analytification}

\noindent
Throughout this section an {\it algebraic variety over $\mathbf C$} means a
reduced separated scheme of finite type over $\mathbf C$, not necessarily
irreducible. Its classical points are its closed points, equivalently its
$\mathbf C$-points. This convention is the scheme-theoretic translation of
the convention in \cite[\S 2, no. 5]{GAGA} and is compatible with the reduced
analytic spaces of Section \ref{section-serre-analytic-spaces}.

\begin{lemma}
\label{lemma-algebraic-charts-analytic}
Let $U \subset \mathbf C^n$ and $U' \subset \mathbf C^{n'}$ be Zariski locally
closed reduced subsets.
\begin{enumerate}
\item The Zariski topology on $\mathbf C^n$ is coarser than its usual topology.
\item The set $U$, with its reduced induced structure, is a Serre complex
analytic space.
\item A regular map $f : U \to U'$ is holomorphic.
\item A biregular isomorphism $U \to U'$ is an analytic isomorphism.
\end{enumerate}
\end{lemma}

\begin{proof}
A Zariski closed subset is the common zero set of polynomials, and polynomials
are continuous and holomorphic. This proves (1) and proves (2) for a closed
subset; a locally closed subset is open in its closure, and restriction gives
the general case. A regular function on a Zariski open subset is locally a
quotient $g/h$ of polynomials with $h$ nonvanishing. Hence it is holomorphic.
Applying this coordinate by coordinate proves (3). Applying (3) to $f$ and to
$f^{-1}$ proves (4).
\end{proof}

\begin{proposition}
\label{proposition-analytification}
Let $X$ be an algebraic variety over $\mathbf C$. There is a unique Serre
complex analytic space $X^{an}$ with the following property: for every
algebraic chart
$$
\varphi : V \longrightarrow U \subset \mathbf C^n
$$
from a Zariski open $V \subset X$ onto a Zariski locally closed reduced
subset, $V$ is open in $X^{an}$ and $\varphi$ is an analytic isomorphism.

The construction is functorial. Every morphism $f : X \to Y$ induces a
holomorphic map $f^{an} : X^{an} \to Y^{an}$, and there is a canonical
morphism of locally ringed spaces
$$
a_X : (X^{an}, \mathcal O_{X^{an}}) \longrightarrow
(X, \mathcal O_X)
$$
whose map on points sends a classical point to the corresponding closed
scheme point and whose map on structure sheaves sends a regular function to
the same function viewed as holomorphic.
\end{proposition}

\begin{proof}
Cover $X$ by finitely many algebraic charts $V_i \to U_i$. Transport the
analytic structure of $U_i$ to $V_i$. On $V_i \cap V_j$ the transition maps
and their inverses are holomorphic by
Lemma \ref{lemma-algebraic-charts-analytic}; the overlap is open for the usual
topology by the same lemma. The cocycle condition is inherited from the
algebraic charts. The sheaves and the locally ringed spaces therefore glue by
Sheaves, Lemma \ref{sheaves-lemma-glue-sheaves-structures}.

It remains to check the Hausdorff axiom, rather than hide it in the word
``glue''. The inverse image of the diagonal of the glued point set in every
$V_i \times V_j$ is the graph of the algebraic identification on the overlap.
Since $X$ is separated, this graph is Zariski closed; it is closed in the usual
topology by Lemma \ref{lemma-algebraic-charts-analytic}. These products form an
open cover, so the diagonal is closed and the glued space is Hausdorff.
Uniqueness follows because the charts cover $X$ and determine both its topology
and its structure sheaf.

For a morphism $f : X \to Y$, its restrictions in algebraic charts are
holomorphic by Lemma \ref{lemma-algebraic-charts-analytic}; they glue to
$f^{an}$, compatibly with identities and composition. The usual topology is
finer than the Zariski topology, so the point map underlying $a_X$ is
continuous. Regular functions are holomorphic in charts, and the induced maps
on stalks are local. This constructs the asserted morphism of locally ringed
spaces.
\end{proof}

\begin{lemma}
\label{lemma-analytification-properties}
Let $X$ and $Y$ be algebraic varieties over $\mathbf C$.
\begin{enumerate}
\item The space $X^{an}$ is locally compact, second countable, and countable
at infinity.
\item There is a canonical analytic isomorphism
$(X \times_{\mathbf C} Y)^{an} \cong X^{an} \times Y^{an}$.
\item If $Z \subset X$ is a Zariski locally closed reduced subvariety, then
$Z^{an}$ is the reduced analytic subspace induced on $Z(\mathbf C)$ by
$X^{an}$.
\end{enumerate}
\end{lemma}

\begin{proof}
A Zariski locally closed subset of $\mathbf C^n$ is locally compact, second
countable, and a countable union of compact subsets. The same properties pass
to a finite union of open chart domains, proving (1). Products of algebraic
charts are algebraic charts and the usual topology on a product of subsets of
complex affine spaces is the product topology; uniqueness in
Proposition \ref{proposition-analytification} proves (2). Restricting charts of
$X$ to $Z$ and using Lemma \ref{lemma-algebraic-charts-analytic} proves (3).
\end{proof}

\section{Local algebra of analytification}
\label{section-local-comparison}

\noindent
Let $x \in X(\mathbf C)$. We write
$$
\theta_x : \mathcal O_{X,x} \longrightarrow \mathcal O_{X^{an},x}
$$
for the local homomorphism induced by $a_X$. Both rings are Noetherian local
rings with residue field $\mathbf C$: this is standard for the algebraic ring,
and for the analytic ring it is Lemma \ref{lemma-local-analytic-algebra}.

\begin{proposition}
\label{proposition-analytification-closed-ideal}
Let $Z \subset X$ be a Zariski locally closed reduced subvariety and let
$x \in Z(\mathbf C)$. After replacing $X$ by a Zariski neighbourhood of $x$,
view $Z$ as closed. If $\mathcal J_{Z,x} \subset \mathcal O_{X,x}$ is its
algebraic ideal and $\mathcal I_{Z^{an},x} \subset
\mathcal O_{X^{an},x}$ is the ideal of holomorphic germs vanishing on
$Z^{an}$, then
$$
\mathcal I_{Z^{an},x} =
\mathcal J_{Z,x}\mathcal O_{X^{an},x}.
$$
Consequently
$$
\mathcal O_{Z^{an},x} \cong
\mathcal O_{X^{an},x}/
\mathcal J_{Z,x}\mathcal O_{X^{an},x}.
$$
\end{proposition}

\begin{proof}
First take $X=\mathbf A^n_{\mathbf C}$ and translate $x$ to the origin. Put
$$
R=\mathbf C[z_1,\ldots,z_n]_{(z_1,\ldots,z_n)},
\qquad
H=\mathbf C\{z_1,\ldots,z_n\}.
$$
Let $J \subset R$ be the radical ideal of $Z$. The local analytic
Nullstellensatz says that the ideal of holomorphic germs vanishing on $Z$ is
$\sqrt{JH}$. We show that $JH$ is already radical.

The Taylor homomorphisms identify the maximal-adic completions of both $R$ and
$H$ with $\mathbf C[[z_1,\ldots,z_n]]$. Completion is exact on finite modules
by Algebra, Lemma \ref{algebra-lemma-completion-tensor}; hence the completion
of $H/JH$ is
$$
\mathbf C[[z_1,\ldots,z_n]]/
J\mathbf C[[z_1,\ldots,z_n]],
$$
which is also the completion of $R/J$. The field $\mathbf C$ is excellent by
More on Algebra, Proposition
\ref{more-algebra-proposition-ubiquity-excellent}, and $R/J$, being a
localization of a finite type $\mathbf C$-algebra, is excellent by More on
Algebra, Lemma \ref{more-algebra-lemma-finite-type-over-excellent}. It is
therefore Nagata by More on Algebra, Lemma
\ref{more-algebra-lemma-quasi-excellent-nagata}. Its
completion is reduced by More on Algebra, Lemma
\ref{more-algebra-lemma-completion-reduced}. Thus the completion of $H/JH$ is
reduced. The ring $H/JH$ injects into its completion by Krull's intersection
theorem, Algebra, Lemma
\ref{algebra-lemma-intersect-powers-ideal-module-zero}; hence $H/JH$ is
reduced and $JH=\sqrt{JH}$. The analytic Nullstellensatz now gives the desired
equality in affine space.

For general $X$, work in an algebraic chart which realizes $X$ and $Z$ as
reduced locally closed subsets of one affine space. Apply the affine result to
their ambient ideals and pass to the quotient. This gives both displayed
formulas.
\end{proof}

\begin{theorem}
\label{theorem-analytification-completed-local-rings}
For every algebraic variety $X$ over $\mathbf C$ and every
$x \in X(\mathbf C)$, the local homomorphism $\theta_x$ induces an isomorphism
of completed local rings
$$
\widehat{\theta}_x :
\widehat{\mathcal O}_{X,x} \xrightarrow{\ \sim\ }
\widehat{\mathcal O}_{X^{an},x}.
$$
In particular, $\theta_x$ is injective.
\end{theorem}

\begin{proof}
For affine space at the origin the map is
$$
\mathbf C[z_1,\ldots,z_n]_{(z_1,\ldots,z_n)}
\longrightarrow \mathbf C\{z_1,\ldots,z_n\},
$$
and both maximal-adic completions are canonically
$\mathbf C[[z_1,\ldots,z_n]]$ by Taylor expansion and truncation. For a
locally closed reduced $X$ in affine space, Proposition
\ref{proposition-analytification-closed-ideal} identifies its analytic local
ring with the quotient by the extended algebraic ideal. Exactness of
completion on finite modules, Algebra, Lemma
\ref{algebra-lemma-completion-tensor}, identifies both completions with the
same quotient of the formal power-series ring. The assertion is local, so
charts prove the general case. Finally, $\mathcal O_{X,x}$ injects into its
completion by Algebra, Lemma
\ref{algebra-lemma-intersect-powers-ideal-module-zero}, which proves the last
sentence.
\end{proof}

\begin{theorem}
\label{theorem-analytification-local-faithfully-flat}
The local homomorphism
$$
\theta_x : \mathcal O_{X,x} \longrightarrow \mathcal O_{X^{an},x}
$$
is faithfully flat. In Serre's terminology,
$(\mathcal O_{X,x},\mathcal O_{X^{an},x})$ is a flat pair: the quotient
$\mathcal O_{X^{an},x}/\mathcal O_{X,x}$ is flat over $\mathcal O_{X,x}$,
and for every ideal $I \subset \mathcal O_{X,x}$ one has
$$
I\mathcal O_{X^{an},x}\cap\mathcal O_{X,x}=I.
$$
\end{theorem}

\begin{proof}
The map on completions is an isomorphism by
Theorem \ref{theorem-analytification-completed-local-rings}, hence is flat.
More on Algebra, Lemma \ref{more-algebra-lemma-flat-completion} implies that
$\theta_x$ is flat, and Algebra, Lemma
\ref{algebra-lemma-local-flat-ff} makes this local flat map faithfully flat.
The contraction formula is Algebra, Lemma
\ref{algebra-lemma-faithfully-flat-universally-injective}. Universal
injectivity after tensoring the exact sequence
$0\to\mathcal O_{X,x}\to\mathcal O_{X^{an},x}$ shows that the displayed
quotient is flat, which is precisely the flat-pair assertion.
\end{proof}

\begin{lemma}
\label{lemma-analytification-local-dimension}
For every $x \in X(\mathbf C)$,
$$
\dim(\mathcal O_{X,x})=\dim(\mathcal O_{X^{an},x}).
$$
\end{lemma}

\begin{proof}
Each ring has the same dimension as its completion by More on Algebra, Lemma
\ref{more-algebra-lemma-completion-dimension}, and the two completions are
isomorphic by Theorem \ref{theorem-analytification-completed-local-rings}.
\end{proof}

\begin{lemma}
\label{lemma-analytification-dimension}
If $X$ is irreducible of dimension $r$, then $X^{an}$ has analytic dimension
$r$ at every point.
\end{lemma}

\begin{proof}
For a closed point $x$, Varieties, Lemma
\ref{varieties-lemma-dimension-locally-algebraic}, parts
\ref{varieties-item-dimension-at-closed-point} and
\ref{varieties-item-dimension-irreducible}, gives
$$
\dim(\mathcal O_{X,x})=r.
$$
Apply
Lemma \ref{lemma-analytification-local-dimension}, followed by
Lemma \ref{lemma-components-dimension-analytic-space}, which identifies
analytic dimension with the Krull dimension of the analytic local ring.
\end{proof}

\section{The classical and Zariski topologies}
\label{section-classical-zariski-topologies}

\begin{lemma}
\label{lemma-zariski-dense-open-analytically-dense}
Let $U \subset X$ be Zariski open and Zariski dense. Then $U^{an}$ is dense
in $X^{an}$.
\end{lemma}

\begin{proof}
Put $Z=X\setminus U$ with its reduced induced structure. If a point
$x \in X^{an}$ were not in the closure of $U^{an}$, then $Z^{an}$ would
contain an analytic neighbourhood of $x$. Its local vanishing ideal at $x$
would therefore be zero. Proposition
\ref{proposition-analytification-closed-ideal} and the injectivity in
Theorem \ref{theorem-analytification-completed-local-rings} would give
$\mathcal J_{Z,x}=0$. Thus $Z$ would agree with $X$ on a nonempty Zariski
neighbourhood of $x$, contradicting the Zariski density of $U$.
\end{proof}

\begin{proposition}
\label{proposition-proper-if-and-only-if-compact-analytification}
An algebraic variety $X$ over $\mathbf C$ is proper over $\mathbf C$ if and
only if the topological space $X^{an}$ is compact.
\end{proposition}

\begin{proof}
Apply Chow's lemma in the precise form of Limits, Lemma
\ref{limits-lemma-chow-finite-type}. After replacing the source by its
reduction as allowed by the following remark there, we obtain a proper
surjective morphism $\pi : U \to X$ and an immersion
$U \to \mathbf P^n_{\mathbf C}$. Let $Y$ be the reduced closure of $U$ in
$\mathbf P^n_{\mathbf C}$. Then $Y$ is projective and $U$ is Zariski open
and dense in $Y$. Moreover the graph of $\pi$ is closed in $X\times Y$:
over $X$ it is the graph of a morphism from the proper $X$-scheme $U$ to the
separated $X$-scheme $X\times Y$.

If $X$ is proper, then $U$ is proper over $\mathbf C$. Its immersion into
projective space is proper and hence closed, so $U=Y$. The space
$Y^{an}$ is a closed subset of the compact Hausdorff space
$\mathbf P^n(\mathbf C)$, hence compact. The continuous surjection
$\pi^{an}:Y^{an}\to X^{an}$ shows that $X^{an}$ is compact.

Conversely, assume $X^{an}$ compact. Then $X^{an}\times Y^{an}$ is compact,
and the analytification of the closed graph is closed in this product by
Lemma \ref{lemma-analytification-properties}. Its projection to $Y^{an}$ is
$U^{an}$, hence $U^{an}$ is compact and therefore closed in the Hausdorff
space $Y^{an}$. It is dense by
Lemma \ref{lemma-zariski-dense-open-analytically-dense}, so $U^{an}=Y^{an}$.
Thus $U=Y$ because a nonempty closed subset of a finite type
$\mathbf C$-scheme has a closed $\mathbf C$-point. Now the projective variety
$Y$ is proper and maps surjectively to $X$; Morphisms, Lemma
\ref{morphisms-lemma-image-proper-is-proper} implies that $X$ is proper.
\end{proof}

\begin{lemma}
\label{lemma-dense-image-contains-open}
Let $f : X \to Y$ be a morphism of algebraic varieties over $\mathbf C$. If
$f(X)$ is Zariski dense in $Y$, then $f(X)$ contains a Zariski open dense
subset of $Y$.
\end{lemma}

\begin{proof}
The image is constructible by Morphisms, Lemma
\ref{morphisms-lemma-chevalley}. A dense constructible subset contains a dense
open subset by Topology, Lemma
\ref{topology-lemma-dense-in-constructible}.
\end{proof}

\begin{proposition}
\label{proposition-image-closure}
Let $f : X \to Y$ be a morphism of algebraic varieties. The closure of
$f(X(\mathbf C))$ in $Y^{an}$ is the set of classical points of the Zariski
closure of $f(X)$ in $Y$.
\end{proposition}

\begin{proof}
Let $T$ be the reduced Zariski closure of $f(X)$. Applying
Lemma \ref{lemma-dense-image-contains-open} to $X\to T$ gives a Zariski open
dense $V\subset T$ contained in $f(X)$. By
Lemma \ref{lemma-zariski-dense-open-analytically-dense}, $V^{an}$ is dense in
$T^{an}$, so $T^{an}$ is contained in the analytic closure of $f(X)$. The
opposite inclusion holds because $T^{an}$ is closed in $Y^{an}$ by
Lemma \ref{lemma-analytification-properties}.
\end{proof}

\section{An analytic criterion for regularity}
\label{section-analytic-regularity}

\begin{proposition}
\label{proposition-algebraic-graph-holomorphic-regular}
Let $X$ and $Y$ be algebraic varieties and let
$f : X^{an}\to Y^{an}$ be holomorphic. If its graph is the classical point
set of a Zariski locally closed reduced subvariety
$T\subset X\times_{\mathbf C}Y$, then $f$ is induced by a morphism
$X\to Y$.
\end{proposition}

\begin{proof}
The first projection $p:T\to X$ is a morphism and is bijective on classical
points. Its inverse $x\mapsto(x,f(x))$ is holomorphic. By
Lemma \ref{lemma-analytification-properties}, $p^{an}$ is therefore an
analytic isomorphism. Proposition
\ref{proposition-analytic-isomorphism-algebraic} below shows that $p$ is an
isomorphism. Hence $f=\operatorname{pr}_Y\circ p^{-1}$ is regular.
\end{proof}

\begin{proposition}
\label{proposition-analytic-isomorphism-algebraic}
Let $p:T\to X$ be a morphism of algebraic varieties which is bijective on
classical points. If $p^{an}:T^{an}\to X^{an}$ is an analytic isomorphism,
then $p$ is an isomorphism.
\end{proposition}

\begin{proof}
First $p$ is a homeomorphism for the Zariski topology on classical points.
Indeed, if $F\subset T$ is Zariski closed, then $F^{an}$ is closed in
$T^{an}$. Its image under the analytic homeomorphism is closed in $X^{an}$.
By Proposition \ref{proposition-image-closure}, applied to $F\to X$, this
image is Zariski closed.

Both schemes are Jacobson by Varieties, Lemma
\ref{varieties-lemma-locally-finite-type-Jacobson}. Hence their closed subsets
are determined by their closed points. Passing from irreducible closed subsets
to their unique generic points (Schemes, Lemma
\ref{schemes-lemma-scheme-sober}) extends the preceding homeomorphism uniquely
to a homeomorphism of the full underlying scheme spaces.

Fix $t\in T(\mathbf C)$ and put $x=p(t)$. The Zariski homeomorphism makes
$$
A=\mathcal O_{X,x}\longrightarrow A'=\mathcal O_{T,t}
$$
injective. The analytic isomorphism identifies
$B=\mathcal O_{X^{an},x}$ with $\mathcal O_{T^{an},t}$. By
Theorem \ref{theorem-analytification-local-faithfully-flat}, both $A$ and
$A'$ embed in this common ring $B$.

The irreducible components through $x$ correspond to those through $t$.
Write their minimal primes as $\mathfrak p_i\subset A$ and
$\mathfrak p'_i\subset A'$, with
$\mathfrak p'_i\cap A=\mathfrak p_i$. Put
$$
K_i=\operatorname{Frac}(A/\mathfrak p_i),
\qquad
K'_i=\operatorname{Frac}(A'/\mathfrak p'_i).
$$
The corresponding irreducible varieties are homeomorphic and have the same
dimension, so $K'_i/K_i$ is a finite extension. To see that its degree is one,
use Morphisms, Lemma \ref{morphisms-lemma-finite-degree} and generic flatness,
Morphisms, Proposition \ref{morphisms-proposition-generic-flatness}, to shrink
the target to a nonempty open over which the component map is finite locally
free. Since the ground field has characteristic zero, $K'_i/K_i$ is separable;
after shrinking once more the map is finite etale of degree
$[K'_i:K_i]$. Every closed fibre then consists of that many reduced
$\mathbf C$-points. Bijectivity of $p$ forces $[K'_i:K_i]=1$.

Let $S$ be the complement in $A$ of the union of its minimal primes and define
$S'$ similarly for $A'$. Lemma \ref{lemma-total-fractions-reduced} below gives
$$
A_S=\prod_iK_i=\prod_iK'_i=A'_{S'}.
$$
If $h\in A'$, write $h=g/s$ in this common total ring of fractions with
$g\in A$ and $s\in S$. Thus $g=sh\in sB$. Faithful flatness and Algebra,
Lemma \ref{algebra-lemma-faithfully-flat-universally-injective}, give
$sB\cap A=sA$, so $g=sa$ for some $a\in A$. Hence $s(h-a)=0$. The element
$s$ lies in no minimal prime of the reduced ring $A'$, and is therefore a
nonzerodivisor by Algebra, Lemma
\ref{algebra-lemma-reduced-ring-sub-product-fields}. Thus $h=a$ and $A'=A$.

We have proved that the local maps are isomorphisms at every closed point.
Thus $p$ is etale at every closed point by Etale, Lemma
\ref{etale-lemma-characterize-etale-Noetherian}. The etale locus is open; its
closed complement, if nonempty, would contain a closed point because $T$ is
Jacobson. Hence $p$ is etale everywhere.

Finally, repeat the finite-degree argument above for the restriction of $p$ to
the reduced closure of an arbitrary point of $T$. The Zariski homeomorphism
identifies this closure with the closure of its image, and the same argument
shows that the induced extension of their function fields, namely the residue
field extension at the two generic points, has degree one. Thus $p$ is
injective on points with trivial residue extensions, hence radicial by
Morphisms, Lemma \ref{morphisms-lemma-universally-injective}. An etale
radicial morphism is an open immersion by Etale, Theorem
\ref{etale-theorem-etale-radicial-open}. Since $p$ is also surjective, it is an
isomorphism.
\end{proof}

\begin{lemma}
\label{lemma-total-fractions-reduced}
Let $A$ be a reduced ring with finitely many minimal primes
$\mathfrak p_1,\ldots,\mathfrak p_r$. Let $S$ be the set of elements lying in
none of the $\mathfrak p_i$ and put
$K_i=\operatorname{Frac}(A/\mathfrak p_i)$. Then
$$
A_S=\prod_{i=1}^rK_i.
$$
\end{lemma}

\begin{proof}
Algebra, Lemma \ref{algebra-lemma-reduced-ring-sub-product-fields} says that
$S$ is exactly the set of nonzerodivisors and that
$A_{\mathfrak p_i}=K_i$. Algebra, Lemma
\ref{algebra-lemma-total-ring-fractions-no-embedded-points} then identifies
the total ring of fractions $A_S$ with the displayed product.
\end{proof}

% External analytic input ledger for this continuation:
% (1) the local analytic Nullstellensatz, used exactly once in the proof of
%     Proposition \ref{proposition-analytification-closed-ideal};
% (2) Taylor expansion and the completion identity
%     \widehat{\mathbf C\{z_1,...,z_n\}}=\mathbf C[[z_1,...,z_n]], used in
%     the two local-comparison proofs;
% (3) compactness and Hausdorffness of complex projective space, used in the
%     proof of Proposition
%     \ref{proposition-proper-if-and-only-if-compact-analytification}.
% Oka--Cartan coherence and Weierstrass Noetherianity were already stated as
% explicit dependencies in the foundation block; no properness-comparison,
% identity, normalization, or unmentioned analytic theorem is invoked here.
